% ---- pg-051 （书页 39） ----
\chapter{矩阵代数与随机向量}
\phantomsection\label{chap:2}
% 章首页 pg-051

\section{引言}  % 原书 2.1
\phantomsection\label{sec:2.1}

在第 1 章中我们看到,多元数据能方便地用一个数字阵列来表示. 一般的,一个 $n$ 行 $p$ 列
的矩形数字阵列称为 $n\times p$ 维矩阵. 矩阵代数的应用将给多元方法的研究带来很多便利.

本章所提供的一些矩阵代数结论将使我们能简明地表述统计模型. 此外,用矩阵表示的
各种关系很容易在计算机上编程,从而使一些重要统计量的计算可由程序来完成.

我们从引入一些极基本的概念开始,这对以后的一些统计方法的几何解释和代数说明是
必不可少的. 如果以前没有接触过基本的矩阵代数知识,可以先学习本章 2.2 节和补充 2A 中
所提供的材料,前者简明扼要,后者比较详细.

\section{矩阵和向量代数基础}  % 原书 2.2
\phantomsection\label{sec:2.2}

\mvsub{向量}

由 $n$ 个实数 $x_1,x_2,\cdots,x_n$ 组成的一个数组 $\bm{x}$ 称为向量,写做

\[
\bm{x}=\begin{bmatrix}x_1\\x_2\\\vdots\\x_n\end{bmatrix}
\quad\text{或}\quad
\bm{x}'=[x_1,x_2,\cdots,x_n]
\]

其中字母右上角的撇($'$)表示将列变换成行的转置运算.

向量 $\bm{x}$ 在几何上可表示为一个 $n$ 维有向线段,它沿第一个轴的分量为 $x_1$,沿第二个轴的
分量为 $x_2,\cdots$,沿第 $n$ 个轴的分量为 $x_n$. 对 $n=3$ 的情形,可用\hyperref[fig:2.1]{图 2.1} 说明.

向量可通过乘一个常数 $c$ 来实现伸长或缩短. 具体地说,
我们定义向量 $c\bm{x}$ 为

\[
c\bm{x}=\begin{bmatrix}cx_1\\cx_2\\\vdots\\cx_n\end{bmatrix}
\]

\begin{mvfigure}
\begin{tikzpicture}[scale=1.1]
\draw[->] (0,0) -- (3.6,0) node[right] {$x_2$};
\draw[->] (0,0) -- (0,2.7) node[above] {$x_1$};
\draw[->] (0,0) -- (-1.3,-0.85) node[below left] {$x_3$};
\draw[dashed] (2.6,0) -- (2.6,2.1) -- (0,2.1);
\draw[dashed] (2.6,2.1) -- (1.5,1.55) (0,0) -- (2.6,2.1);
\draw[dashed] (1.5,1.55) -- (1.5,0.9) -- (2.6,2.1);
\draw[dashed] (1.5,0.9) -- (-0.5,0.1) (1.5,1.55) -- (-0.5,0.1);
\draw[very thick,->] (0,0) -- (1.5,1.55);
\draw (-0.05,2.1) -- (0.05,2.1) node[left,pos=0] {\small 2};
\draw (2.55,0.05) -- (2.55,-0.05) node[below] {\small 3};
\draw (0,0.95) node[left,inner sep=1pt] {\small 0};
\draw (1.75,0.65) node {\small 1};
\draw (0.62,0.95) node {\small $\bm{x}$};
\end{tikzpicture}
\phantomsection\label{fig:2.1}\mvcap{图 2.1\quad 向量 $\bm{x}'=[1,3,2]$}
\end{mvfigure}

即 $c\bm{x}$ 是用 $c$ 去乘 $\bm{x}$ 的每个元素而得到的向量[见\hyperref[fig:2.2]{图 2.2}(a)].

两个向量可以相加. $\bm{x}$ 和 $\bm{y}$ 的加法定义为
